\section{单圈阶的等价性}
\label{sec:equiv}

正如第~\ref{sec:ope}~节所证明的，准 PDF 与赝 PDF 及其匹配系数通过式~(\ref{eq:equiv})与式~(\ref{eq:ftc-C})中的简单傅里叶变换相联系。此关系对微扰论所有阶都成立。本节在单圈阶检验式~(\ref{eq:equiv})与式~(\ref{eq:ftc-C})。我们的主要展示取 $\Gamma=\gamma^0$，同时也给出 $\Gamma=\gamma^z$ 情形的最后结果。

在费曼规范下，我们在维数正规化 $d=4-2\epsilon$ 中单圈计算非极化同位旋矢量准 PDF、赝 PDF 与光锥 PDF 的夸克矩阵元。外夸克态取在壳且无质量，紫外发散与共线发散分别用 $1/\epsilon_{\mbox{\tiny UV}}\ (\epsilon_{\mbox{\tiny UV}}>0)$ 与 $1/\epsilon_{\mbox{\tiny IR}}\ (\epsilon_{\mbox{\tiny IR}}<0)$ 正则化。单圈费曼图示于图~\ref{fig:diagram}。


\begin{widetext}

\begin{figure*}[t!]
	\centering
	\includegraphics[width=.9\textwidth]{images/quasi_diagram.pdf}
	\caption{准 PDF、\spcorr 与赝 PDF 的单圈费曼图。第一张称为``顶点图''（vertex），第二、三张称为``帆图''（sail），最后一张称为``蝌蚪图''（tadpole）。图中还包含标准的夸克自能波函数图。}
	\label{fig:diagram}
\end{figure*}

在动量为 $p^\mu=(p^0=p^z,0,0,p^z)$ 的在壳夸克态中，对 $\Gamma=\gamma^0$，各图给出
\begin{align}  \label{eq:qtloopint1}
\tilde{Q}^{(1)}_\text{vertex}(\zeta,z^2,\epsilon)
&= {\mu^{2\epsilon} \iota^\epsilon \over 2p^0}\bar{u}(p)\int {d^d k\over (2\pi)^d}(-ig T^a\gamma^\mu) {i\over \slashed k} \gamma^0 {i\over\slashed k} (-i gT^a\gamma^\nu){ - i g_{\mu\nu} \over (p-k)^2}u(p)e^{-i k^z z}\ ,
  \\
\tilde{Q}^{(1)}_\text{sail}(\zeta,z^2,\epsilon)
&= {\mu^{2\epsilon} \iota^\epsilon\over 2p^0}\bar{u}(p)\int {d^d k\over (2\pi)^d} (ig T^a \gamma^0) {1\over i(p^z-k^z)} \left( e^{-ip^z z}- e^{-ik^z z}\right)  \delta^{\mu z} {i\over \slashed k} (-ig T^a \gamma^\nu) {-ig_{\mu\nu}\over (p-k)^2}u(p)
  \nonumber\\
&+ {\mu^{2\epsilon}\iota^\epsilon \over 2p^0}\bar{u}(p) \int {d^d k\over (2\pi)^d} (-ig T^a \gamma^\nu){i\over \slashed k} (ig T^a \gamma^0) {1\over i(p^z-k^z)} \left( e^{-ip^z z}- e^{-ik^z z}\right)  \delta^{\mu z}   {-ig_{\mu\nu}\over (p-k)^2}u(p)\ ,
  \nn \\
\tilde{Q}^{(1)}_\text{tadpole}(\zeta,z^2,\epsilon)
&= {\mu^{2\epsilon}\iota^\epsilon \over 2p^0}\bar{u}(p)\int {d^d k\over (2\pi)^d} (-g^2)C_F \gamma^0 \delta^{\mu z}\delta^{\nu z}\left({e^{-ip^z z}- e^{-ik^z z}\over (p^z-k^z)^2} - {ze^{-ip^z z}\over i(p^z-k^z)}\right){-ig_{\mu\nu}\over (p-k)^2}u(p)\ ,
  \nn
\end{align}
其中引入 $\iota = e^{\gamma_E}/(4\pi)$ 是为了在 $\overline{\rm MS}$ 方案中实现 $\mu$，$C_F=4/3$，而 $T^a$ 是基本表示的 SU(3) 色矩阵。末行方括号中正比于 $z$ 的第二项对圈积分没有贡献，因为它在交换 $p^z-k^z\to -(p^z-k^z)$ 下为奇。夸克自能修正为 $\tilde{Q}_{\rm w.fn.}^{(1)}(\zeta,z^2,\epsilon) =\delta Z_\psi\: \tilde{Q}^{(0)}(\zeta,z^2)$，其中树图矩阵元 $\tilde{Q}^{(0)}(\zeta,z^2)=e^{-i\zeta}$，而在壳重正化常数 $\delta Z_\psi$ 为
\beq
\delta Z_\psi = {\alpha_sC_F\over 2\pi}\left(-{1\over2}\right)\left({1\over\epsilon_{\mbox{\tiny UV}}} - {1\over \epsilon_{\mbox{\tiny IR}}}\right)\ .
\eeq

按照文献~\cite{Ji:2017rah}的方法完成式~(\ref{eq:qtloopint1})中的圈积分后，我们得到
\begin{align}\label{eq:1loopdiagram}
\tilde{Q}^{(1)}_{\rm vertex}(\zeta,z^2,\epsilon)
=\,& {\alpha_sC_F\over 2\pi}e^{\epsilon\gamma_E} \int_0^1 du\ (1-\epsilon)(1-u) e^{-iu \zeta}\Gamma(-\epsilon) 4^{-\epsilon} \big(\mu |z|\big)^{2\epsilon}
  \\
 =\,& {\alpha_sC_F\over 2\pi} e^{\epsilon\gamma_E}\biggl({\mu|z|\over2}\biggr)^{2\epsilon}
 \frac{(-1)\Gamma(2-\epsilon)}{\epsilon_{\mbox{\tiny IR}}}
 {1-i\zeta - e^{-i\zeta}\over \zeta^2}
  \,, \nonumber\\
\tilde{Q}^{(1)}_{\rm sail}(\zeta,z^2,\epsilon)
=\, & {\alpha_sC_F\over 2\pi}  e^{\epsilon\gamma_E}\left[(i\zeta) \int_0^1 du \int_0^1 dt (2-u) e^{-i(1-ut)\zeta} \Gamma(-\epsilon) 4^{-\epsilon}(t^2z^2\mu^2)^{\epsilon}\right.
  \nonumber\\
 &\left. -\int_0^1 du \ e^{-iu \zeta}\Gamma(-\epsilon) 4^{-\epsilon}(z^2\mu^2)^{\epsilon}  + \left({1\over\epsilon_{\mbox{\tiny UV}}} - {1\over \epsilon_{\mbox{\tiny IR}}}\right) e^{-i\zeta}\right]
  \nonumber\\
 =\,& {\alpha_sC_F\over 2\pi} e^{\epsilon\gamma_E} \left({\mu|z|\over2}\right)^{2\epsilon}
 \Bigg\{\frac{-\Gamma(1-\epsilon)}{\epsilon_{\mbox{\tiny IR}}}
 {2(1-\epsilon)\over 1-2\epsilon} \left[{1-e^{-i\zeta}\over -i\zeta} + e^{-i\zeta}(-i\zeta)^{-2\epsilon}\big(\Gamma(2\epsilon)-\Gamma(2\epsilon,-i\zeta)\big)
 \right]
  \nn\\
 & + {\Gamma(1-\epsilon)\over \epsilon_{\mbox{\tiny IR}} } \frac{e^{-i\zeta}}{\epsilon} + \left({1\over\epsilon_{\mbox{\tiny UV}}} - {1\over \epsilon_{\mbox{\tiny IR}}}\right) e^{-i\zeta} \Bigg\}
  \,, \nn\\
\tilde{Q}^{(1)}_{\rm tadpole}(\zeta,z^2,\epsilon)
=\, & {\alpha_sC_F\over 2\pi} e^{\epsilon\gamma_E}\left({\mu|z|\over2}\right)^{2\epsilon}  {\Gamma(1-\epsilon)\over \epsilon_{\mbox{\tiny UV}} (1-2\epsilon)}e^{-i\zeta}
 \ . \nn
\end{align}
为简单起见，引用式~(\ref{eq:1loopdiagram})中的单圈结果时我们略去了树图的乘法旋量因子 $\bar u\gamma^0 u$，下文引用\spcorr、准 PDF 与赝 PDF 结果时也将如此。
由于 $\alpha_s^{\rm bare}=\alpha_s(\mu) \mu^{2\epsilon} Z_g^2 = \alpha_s(\mu) \mu^{2\epsilon} + {\cal O}(\alpha_s^2)$ 与 $\mu$ 无关，我们在裸函数中不把 $\mu$ 写为宗量。
在每个结果的最终表达式中，我们还标明了 $1/\epsilon$ 因子（$\epsilon\to 0$ 展开后保留者）是红外还是紫外发散。结合波函数修正，裸\spcorr $\tilde{Q}^{(1)}(\zeta,z^2,\epsilon)$ 为
\begin{align} \label{eq:1loopioffe}
\tilde{Q}^{(1)}(\zeta,z^2,\epsilon)
&= {\alpha_sC_F\over 2\pi}e^{\epsilon\gamma_E} \left({\mu|z|\over2}\right)^{2\epsilon} \Bigg\{
 {3\over2}\left({1\over \epsilon_{\mbox{\tiny UV}}}-{1\over \epsilon_{\mbox{\tiny IR}}}\right)e^{-i\zeta}
 +  \frac{(-1)\Gamma(2-\epsilon)}{\epsilon_{\mbox{\tiny IR}}}
  \Biggl[ {(1-i\zeta - e^{-i\zeta})\over \zeta^2}
 \\
& \qquad\qquad\qquad\qquad\qquad\quad 
 +{2\over 1-2\epsilon}\left\{ {1-e^{-i\zeta}\over -i\zeta} + e^{-i\zeta}(-i\zeta)^{-2\epsilon}
 \big(\Gamma(2\epsilon)-\Gamma(2\epsilon,-i\zeta)\big) - \frac{e^{-i\zeta}}{2\epsilon} \right\} \Biggr]
 \Bigg\}\ . \nn
\end{align}
注意当 $\epsilon\to 0$ 时最内层花括号中的项不含 $1/\epsilon$ 项。我们还可以验证，在算符的定域极限下，正如矢量流守恒所预期的那样，裸单圈修正为零：
\begin{align} \label{eq:Qvcc}
\lim_{z\to 0} \tilde{Q}^{(1)}(z P^z,z^2,\epsilon)
  = \lim_{z\to 0}  {\alpha_sC_F\over 2\pi}e^{\epsilon\gamma_E} \left({\mu|z|\over2}\right)^{2\epsilon} \bigg\{
 {3\over2}\left({1\over \epsilon_{\mbox{\tiny UV}}}-{1\over \epsilon_{\mbox{\tiny IR}}}\right)
  + \frac{3}{2\epsilon_{\mbox{\tiny IR}}} \bigg\} = 0
 \,,
\end{align}
这里我们注意到 $1/\epsilon_{\mbox{\tiny IR}}$ 项的相消很重要，因为 $\epsilon>0$ 的假设只对 $1/\epsilon_{\mbox{\tiny UV}}$ 项成立。
相应的裸赝 PDF 为
\begin{align}\label{eq:1looppseudo}
{\cal P}^{(1)}(x,z^2,\epsilon)
 &= {\alpha_sC_F\over 2\pi} \left[-1+x + \frac{2}{(1-2\epsilon)} -
 \left( {2\over (1-2\epsilon)}{1\over (1-x)^{1-2\epsilon}}\right)^{[0,1]}_{+(1)}\right]e^{\epsilon\gamma_E}\left({\mu|z|\over2}\right)^{2\epsilon} \frac{\Gamma(2-\epsilon)}{\epsilon_{\mbox{\tiny IR}}} \theta(x)\theta(1-x)\\
 &\quad +{\alpha_sC_F\over 2\pi} \:
  e^{\epsilon\gamma_E}\left({\mu|z|\over2}\right)^{2\epsilon}{3\over2}\left({1\over \epsilon_{\mbox{\tiny UV}}}-{1\over \epsilon_{\mbox{\tiny IR}}}\right)\, \delta(1-x)
 \ .\nn
\end{align}
由于下面会遇到不同定义域上的 plus 函数，我们定义给定区域 $D$ 内 $x=x_0$ 处的 plus 函数使得
\beq
	\int_D dx\ \big[ g(x)\big]_{+(x_0)}^D\, h(x) = \int_D dx\ g(x) \left[ h(x) - h(x_0)\right]\,.
\eeq
（更多细节见附录~\ref{sec:ep}。）
容易确认，裸赝 PDF 满足定域矢量流守恒 $\lim_{z\to 0} \int\! dx \, {\cal P}^{(1)}(x,z^2\mu^2,\epsilon) =0$，其相消方式与式~(\ref{eq:Qvcc})相同。

现在，根据准 PDF 与\spcorr 或赝 PDF 之间的关系（式~(\ref{eq:fac-tq-orig0},\ref{eq:equiv})），我们可以对式~(\ref{eq:1loopioffe},\ref{eq:1looppseudo})的结果作一次或双重傅里叶变换得到准 PDF。尽管直接了当，傅里叶变换却很微妙，细节见附录~\ref{sec:ft}。这里我们直接引用裸准 PDF 的结果：
\begin{align}  \label{eq:1loopquasi}
\tilde{q}^{(1)}(x,p^z,\epsilon)
	=& {\alpha_sC_F\over 2\pi}\left\{{3\over2}\left({1\over \epsilon_{\mbox{\tiny UV}}}-{1\over \epsilon_{\mbox{\tiny IR}}}\right)\delta(1-x) + {\Gamma(\epsilon+{1\over2})e^{\epsilon\gamma_E}\over \sqrt{\pi}}{\mu^{2\epsilon}\over p_z^{2\epsilon}}{1-\epsilon \over \epsilon_{\mbox{\tiny IR}} (1-2\epsilon)} \right.\\
	&\times\left.\left[|x|^{-1-2\epsilon}\left(1+x+{x\over2}(x-1+2\epsilon)\right)-  |1-x|^{-1-2\epsilon}\left(x+{1\over2}(1-x)^2\right) + I_3(x)\right] \right\}\,,\nn
\end{align}
其中
\beq
	I_3(x) = \theta(x-1)\left(x^{-1-2\epsilon}\over x-1\right)^{[1,\infty]}_{+(1)} -\theta(x)\theta(1-x) \left(x^{-1-2\epsilon}\over 1-x\right)^{[0,1]}_{+(1)} - \delta(1-x) \pi\csc(2\pi\epsilon) + \theta(-x) {|x|^{-1-2\epsilon}\over x-1}\,.
\eeq
经过一些代数运算可以确认，裸准 PDF 满足定域矢量流守恒，即 $\int\! dx\, \tilde{q}^{(1)}(x,p^z,\epsilon) = 0$。为验证这一结果，必须小心区分两类 $1/\epsilon$ 因子：要求 $\epsilon>0$ 以获得 $x=\pm \infty$ 处收敛所产生的 $1/\epsilon_{\mbox{\tiny UV}}$ 因子，以及要求 $\epsilon<0$ 以获得 $x=1$ 处收敛所产生的 $1/\epsilon_{\mbox{\tiny IR}}$ 因子。

得到准 PDF 的另一种方法是直接从费曼图计算它：先对被积函数把 $z$ 傅里叶变换为 $xp^z$。这样，因子 $(e^{-ip^z z}- e^{-ik^z z})$ 变换为 $[\delta(p^z-xp^z)-\delta(k^z-xp^z)]$，而所有圈积分化为 $(d-1)$ 维。这是文献~\cite{Xiong:2013bka,Stewart:2017tvs,Wang:2017qyg}中准 PDF 匹配计算所用的步骤，它与式~(\ref{eq:1loopdiagram}--\ref{eq:1loopquasi})先完全积完积分、之后再作傅里叶变换的做法不同。作为交叉检验，我们在附录~\ref{sec:quasi}中确认了两种步骤得到完全相同的裸准 PDF 式~(\ref{eq:1loopquasi})。

现在考虑 $\epsilon$ 展开，以得到\spcorr、赝 PDF 与准 PDF 的 $\overline{\rm MS}$ 重正化结果。把\spcorr 按 $\epsilon$ 展开得
\begin{align}
 \tilde{Q}^{(1)}(\zeta,z^2,\epsilon)=\delta \tilde{Q}^{(1)}(\zeta,z^2,\mu,\epsilon_{\mbox{\tiny UV}})+\tilde{Q}^{(1)}(\zeta,z^2,\mu,\epsilon_{\mbox{\tiny IR}})+{\cal O}(\epsilon)
\end{align}
其 $\overline{\rm MS}$ 抵消项与重正化\spcorr 为
\begin{align} \label{eq:1loopiofferen}
\delta \tilde{Q}^{(1)}(\zeta,z^2,\mu,\epsilon_{\mbox{\tiny UV}})
&= {\alpha_sC_F\over 2\pi}
  e^{-i\zeta} {3\over2}  {1\over \epsilon_{\mbox{\tiny UV}}}
 \,, \\
\tilde{Q}^{(1)}(\zeta,z^2,\mu,\epsilon_{\mbox{\tiny IR}})
&= {\alpha_sC_F\over 2\pi} \Bigg\{
  \frac{3}{2} \Big(\ln{\mu^2 z^2 e^{2\gamma_E} \over 4} +1\Big) e^{-i\zeta}
  + \Bigl( -\frac{1}{ \epsilon_{\mbox{\tiny IR}} } - \ln\frac{z^2\mu^2e^{2\gamma_E}}{4} - 1 \Bigr) h(\zeta)
  + \frac{2(1 \!-\!i\zeta\!-\! e^{-i\zeta})}{\zeta^2}
  \nn\\
&\qquad\qquad\quad
  + 4 i \zeta e^{-i\zeta}\, {}_3F_3(1,1,1,2,2,2,i\zeta)
  \Bigg\}
  \,. \nn
\end{align}
这里 ${}_3F_3$ 是超几何函数，$\big[(1+x^2)/(1-x)\big]_{+(1)}^{[0,1]}$ 的傅里叶变换给出的函数为
\begin{align}
  h(\zeta) &=\frac{3}{2}\, e^{-i\zeta}
  + \frac{1+i\zeta-e^{-i\zeta}-2i\zeta e^{-i\zeta}}{\zeta^2}
  - 2 e^{-i\zeta} \big[ \Gamma(0,-i\zeta)\!+\!\gamma_E\!+\!\ln(-i\zeta) \big]
  \,.
\end{align}
对坐标空间 PDF 有
\begin{align}
  Q^{(1)}(\zeta,\mu,\epsilon_{\mbox{\tiny IR}})
  &= -  {\alpha_sC_F\over 2\pi}  \frac{1}{  \epsilon_{\mbox{\tiny IR}} }  h(\zeta) \,.
\end{align}

接下来把式~(\ref{eq:1looppseudo})的裸赝 PDF 按 $\epsilon$ 展开，得到 $\overline{\rm MS}$ 抵消项与重正化赝 PDF ${\cal P}^{(1)} (x,z^2,\epsilon)= \delta {\cal P}^{(1)} (x,z^2\mu^2,\epsilon_{\mbox{\tiny UV}})+{\cal P}^{(1)} (x,z^2\mu^2,\epsilon_{\mbox{\tiny IR}})+{\cal O}(\epsilon)$：
\begin{align}\label{eq:1looppseudopdfb}
\delta {\cal P}^{(1)} &(x,z^2\mu^2,\epsilon_{\mbox{\tiny UV}})
	= {\alpha_sC_F\over 2\pi}\:
    {3\over 2\epsilon_{\mbox{\tiny UV}}} \: \delta(1-x)
  \,, \\
{\cal P}^{(1)} &(x,z^2\mu^2,\epsilon_{\mbox{\tiny IR}})
\nn\\
	=&{\alpha_sC_F\over 2\pi}\left\{\left({1+x^2\over 1-x}\right)^{[0,1]}_{+(1)}\left[- \left({1\over \epsilon_{\mbox{\tiny IR}}}+\ln{e^{2\gamma_E}\over 4}\right)-\ln(z^2 \mu^2)-1\right] - \left(4\ln(1-x)\over 1-x\right)^{[0,1]}_{+(1)}  +2(1-x)\right\}\theta(x)\theta(1-x)\nonumber\\
	&  +{\alpha_sC_F\over 2\pi}\left[{3\over2}\ln(z^2\mu^2) + {3\over2}\ln{e^{2\gamma_E}\over 4}+{3\over2}\right]\delta(1-x)
 \,.\nn
\end{align}
注意重正化的 $\overline{\rm MS}$ 赝 PDF 显式地依赖 $\mu^2$，并满足与重正化 $\overline{\rm MS}$ \spcorr 的关系式~(\ref{eq:pseudoRen})。
还有一点很有意思：按 $\epsilon$ 展开之后，重正化 $\overline{\rm MS}$ 赝 PDF 的极限不再满足定域矢量流守恒，因为
\begin{align} \label{eq:Ppdfconsrv}
\lim_{z\to 0} \int\! dx \, {\cal P}^{(1)} (x,z^2\mu^2,\epsilon_{\mbox{\tiny IR}}) \simeq (3\alpha_s C_F/4\pi) \lim_{z\to 0} \ln(z^2\mu^2)
\end{align}
给出发散的结果。同样的发散也出现在单圈 $\overline{\rm MS}$ 重正化\spcorr 式~(\ref{eq:1loopiofferen})中。虽然在 $\overline{\rm MS}$ 中如此，但在其他重正化方案中未必如此。

对准 PDF 而言，重正化计算有两种方法可以考虑：或者展开式~(\ref{eq:1loopquasi})的裸结果并在 $(x,p^z)$ 空间重正化，或者遵循我们偏好的定义式~(\ref{eq:fac-tq-orig0})、把式~(\ref{eq:1loopiofferen})的重正化\spcorr 作傅里叶变换。虽然这两种做法在实际应用中将导致相同的最终 $C$ 结果，但存在一个我们将予以解释的微妙差别。

先考虑在 $(x,p^z)$ 空间进行的准 PDF 重正化。把式~(\ref{eq:1loopquasi})按 $\epsilon$ 展开，并写成 $\tilde{q}^{(1)}(x,p^z,\epsilon) =\delta \tilde{q}^{(1)}(x,\mu/|p^z|,\epsilon_{\mbox{\tiny UV}})+\tilde{q}^{(1)}(x,\mu/|p^z|,\epsilon_{\mbox{\tiny IR}}) +{\cal O}(\epsilon)$，便可识别出 $\overline{\rm MS}$ 抵消项与重正化准 PDF：
\begin{align}\label{eq:qPDFren}
\delta \tilde{q}^{(1)}(x,\mu/|p^z|,\epsilon_{\mbox{\tiny UV}})
=&\,  {\alpha_sC_F\over 2\pi} {3\over 2\epsilon_{\mbox{\tiny UV}}}
 \left[\delta(1-x) - {1\over2}{1\over x^2}\delta^+\Big({1\over x}\Big) - {1\over2}{1\over (1-x)^2}\delta^+\Big({1\over 1-x}\Big) \right]
 \,, \\
\tilde{q}^{(1)}(x,\mu/|p^z|,\epsilon_{\mbox{\tiny IR}})
=&\, {\alpha_sC_F\over 2\pi}\left\{
\begin{array}{ll}
\displaystyle \left({1+x^2\over 1-x}\ln {x\over x-1} + 1 + {3\over 2x}\right)^{[1,\infty]}_{+(1)}- \left({3\over 2x}\right)^{[1,\infty]}_{+(\infty)} &\, x>1\\[10pt]
\displaystyle \left({1+x^2\over 1-x}\left[- {1\over \epsilon_{\mbox{\tiny IR}}} - \ln{\mu^2\over 4p_z^2} + \ln\big(x(1-x)\big)\right] - {x(1+x)\over 1-x}\right)^{[0,1]}_{+(1)}  &\, 0<x<1\\[10pt]
\displaystyle  \left(-{1+x^2\over 1-x}\ln {-x\over 1-x} - 1 + {3\over 2(1-x)}\right)^{[-\infty,0]}_{+(1)} - \left({3\over 2(1-x)}\right)^{[-\infty,0]}_{+(-\infty)} \quad &\, x<0
\end{array}\right.\nn\\[5pt]
& + {\alpha_sC_F\over 2\pi}\left[\delta(1-x) - {1\over2}{1\over x^2}\delta^+\Big({1\over x}\Big) - {1\over2}{1\over (1-x)^2}\delta^+\Big({1\over 1-x}\Big) \right] \left( {3\over2}\ln{\mu^2\over 4p_z^2} + {5\over2}\right)
 \,,\nn
\end{align}
式~(\ref{eq:1loopquasi})的 $\epsilon$ 展开的细节见附录~\ref{sec:ep}，包括此处结果中出现的 plus 函数以及 $x_0=\pm \infty$ 处 $\delta$ 函数的定义。由式~(\ref{eq:qPDFren})所得的 $\overline{\rm MS}$ 准 PDF 仍满足矢量流守恒
\begin{align} \label{eq:qpdfconsrv}
 \int\! dx\: \tilde{q}^{(1)}(x,\mu/|p^z|,\epsilon_{\mbox{\tiny IR}})  =0 \,.
\end{align}
这对各自积分为零的 plus 函数项显然成立，对式~(\ref{eq:qPDFren})中出现的 $\delta$ 函数组合也成立。

式~(\ref{eq:qPDFren})中的重正化 $\overline{\rm MS}$ 准 PDF 与用我们的定义式~(\ref{eq:fac-tq-orig0})所得者略有不同。改用式~(\ref{eq:fac-tq-orig0})以及式~(\ref{eq:1loopiofferen})的重正化\spcorr，我们得到
\begin{align}\label{eq:qPDFft}
\delta \tilde{q}'^{(1)}(x,\mu/|p^z|,\epsilon_{\mbox{\tiny UV}})
=&\,  {\alpha_sC_F\over 2\pi} {3\over 2\epsilon_{\mbox{\tiny UV}}}\delta(1-x)\,, \\
\tilde{q}'^{(1)}(x,\mu/|p^z|,\epsilon_{\mbox{\tiny IR}})
=&\, {\alpha_sC_F\over 2\pi}\left\{
\begin{array}{ll}
\displaystyle \left({1+x^2\over 1-x}\ln {x\over x-1} + 1 + {3\over 2x}\right)^{[1,\infty]}_{+(1)}- \left({3\over 2x}\right)^{[1,\infty]}_{+(\infty)} &\, x>1\\[10pt]
\displaystyle \left({1+x^2\over 1-x}\left[- {1\over \epsilon_{\mbox{\tiny IR}}} - \ln{\mu^2\over 4p_z^2} + \ln\big(x(1-x)\big)\right] - {x(1+x)\over 1-x}\right)^{[0,1]}_{+(1)}  &\, 0<x<1\\[10pt]
\displaystyle  \left(-{1+x^2\over 1-x}\ln {-x\over 1-x} - 1 + {3\over 2(1-x)}\right)^{[-\infty,0]}_{+(1)} - \left({3\over 2(1-x)}\right)^{[-\infty,0]}_{+(-\infty)} \quad &\, x<0
\end{array}\right.\nn\\[5pt]
& + {\alpha_sC_F\over 2\pi}\left[\delta(1-x)\left( {3\over2}\ln{\mu^2\over 4p_z^2} + {5\over2}\right) +{3\over2}\gamma_E\left({1\over(x-1)^2}\delta^+({1\over x-1}) + {1\over(1-x)^2}\delta^+({1\over 1-x})\right)\right]
 \nn .
\end{align}
为完成这一计算，我们把奇异函数 $\ln(\zeta^2)$ 的傅里叶变换定义为
\begin{align} \label{eq:subtlety}
\int {d\zeta\over 2\pi}\ e^{i x \zeta} \ln\zeta^2
= & \biggl[ \left.{d\over d\eta} \int {d\zeta\over 2\pi}\ e^{i x \zeta} (\zeta^2)^\eta \biggr] \right|_{\eta=0}
=\biggl[ \left.{d\over d\eta} {4^\eta\over \Gamma(-\eta)} {\Gamma(\eta+1/2)\over \sqrt{\pi}} {[\theta(x)+\theta(-x)]\over |x|^{1+2\eta}} \biggr]
\right|_{\eta=0}\\
=& \gamma_E \left[\left(-\delta(x) + {1\over x^2} \delta^+\Big({1\over x}\Big)\right) + \left(-\delta(x) + {1\over x^2} \delta^+\Big({1\over -x}\Big)\right)\right] \nn\\
& -\left[\left({1\over x}\right)_{+(0)}^{[0,1]} + \left({1\over x}\right)_{+(\infty)}^{[1,\infty]}\right]\theta(x) - \left[\left({1\over -x}\right)_{+(0)}^{[-1,0]} + \left({1\over -x}\right)_{+(\infty)}^{[-\infty,-1]}\right]\theta(-x)\,,\nn
\end{align}
其中推导第二个等号与最后一个等号时用了式~(\ref{eq:ft},\ref{eq:expansion})的结果，并在需要时取 $\eta\to 0^+$ 或 $\eta\to 0^-$ 极限。这个 $\tilde{q}'^{(1)}(x,\mu/|p^z|,\epsilon_{\mbox{\tiny IR}})$ 不满足矢量流守恒，它与式~(\ref{eq:qPDFren})中的 $\tilde{q}^{(1)}(x,\mu/|p^z|,\epsilon_{\mbox{\tiny IR}})$ 只差 $x_0=\pm\infty$ 处的 $\delta$ 函数。在函数域 $-\infty < x < \infty$ 内二者完全相同。下面我们将看到，式~(\ref{eq:qPDFren})与式~(\ref{eq:qPDFft})最终导致同一个单圈匹配系数结果。


匹配计算所需的最后一个要素是 PDF，其单圈裸矩阵元可以写成 $\overline{\rm MS}$ 抵消项与重正化矩阵元之和 $q^{(1)}(x,\epsilon) = \delta q^{(1)}(x,\epsilon_{\mbox{\tiny UV}})+q^{(1)}(x,\epsilon_{\mbox{\tiny IR}})$：
\begin{align}
 \delta q^{(1)}(x,\epsilon_{\mbox{\tiny UV}}) &= 	{\alpha_sC_F\over 2\pi} {1\over \epsilon_{\mbox{\tiny UV}}} \left({1+x^2\over 1-x}\right)^{[0,1]}_{+(1)}  \theta(x)\theta(1-x)
  \,, \\
 q^{(1)}(x,\epsilon_{\mbox{\tiny IR}}) &= 	{\alpha_sC_F\over 2\pi} {(-1)\over \epsilon_{\mbox{\tiny IR}}} \left({1+x^2\over 1-x}\right)^{[0,1]}_{+(1)} \theta(x)\theta(1-x)
  \,. \nn
\end{align}

有了以上结果，我们现在可以确定到单圈阶的匹配系数。利用式~(\ref{eq:ps-q-fact})可得
\beq \label{eq:pseudoC}
{\cal C}^{(1)}(\alpha,z^2\mu^2)= {\cal P}^{(1)}(\alpha,z^2\mu^2,\epsilon_{\mbox{\tiny IR}}) - q^{(1)}(\alpha,\epsilon_{\mbox{\tiny IR}}) \,.
\eeq
因此，联系 $\Gamma=\gamma^0$ 时赝 PDF 与 PDF 的 $\overline{\rm MS}$ 方案匹配系数为\footnote{\spcorr 在坐标空间的单圈分析最近也见于文献~\cite{Radyushkin:2017lvu,Radyushkin:2018cvn}。我们关于\spcorr 的因子化结果式~(\ref{eq:io-Q-fact})与文献~\cite{Radyushkin:2017lvu}式 (3.35) 以及文献~\cite{Radyushkin:2018cvn}式 (17) 中所给简化\spcorr 的硬部分形式相近。因此把我们的 ${\cal C}(\alpha,z^2\mu^2)/C_0(\mu^2z^2)$ 与该硬部分作比较是有趣的。我们的 $\overline{\rm MS}$ 结果式~(\ref{eq:ps-c})因出现 $2(1-\alpha)$ 项而不同于文献~\cite{Radyushkin:2017lvu}。文献~\cite{Radyushkin:2018cvn}最终版本的结果与我们一致。式~(\ref{eq:ps-c})也与文献~\cite{Ji:2017rah}最初导出的结果一致，只是加上了我们的 $e^{2\gamma_E}$ 项。应当使用式~(\ref{eq:ps-c})中的 ${\cal C}(\alpha,z^2\mu^2)$ 从 MSbar 的赝 PDF 结果提取 MSbar 的 PDF。}
\begin{align} \label{eq:ps-c}
{\cal C}(\alpha,z^2\mu^2)
	=& \left[ 1+ {\alpha_sC_F\over 2\pi}\left({3\over2}\ln(z^2\mu^2)+{3\over2}\ln{e^{2\gamma_E}\over 4}+{3\over2}\right)\right]\delta(1-\alpha)
 \\
 	&+ {\alpha_sC_F\over 2\pi}\left\{\left({1+\alpha^2\over 1-\alpha}\right)^{[0,1]}_{+(1)}\left[-\ln(z^2 \mu^2)-\ln{e^{2\gamma_E}\over 4}-1\right] - \left(4\ln(1-\alpha)\over 1-\alpha\right)^{[0,1]}_{+(1)}  +2(1-\alpha)\right\}\theta(\alpha)\theta(1-\alpha)
  \,. \nn
\end{align}
如预期那样，该结果与红外调节无关。
我们还计算了 $\Gamma=\gamma^z$ 情形的匹配系数，它为 ${\cal C}_{\gamma^z}(\alpha,z^2\mu^2) = {\cal C}(\alpha,z^2\mu^2) + \Delta {\cal C}_{\gamma^z}(\alpha,z^2\mu^2)$，其中
\beq
\Delta {\cal C}_{\gamma^z}(\alpha,z^2\mu^2) = {\alpha_sC_F\over 2\pi} 2(1-\alpha)\theta(\alpha)\theta(1-\alpha)\,.
\eeq
由于式~(\ref{eq:ps-c})中 $\ln(z^2\mu^2)\delta(1-\alpha)$ 项的存在，$\overline{\rm MS}$ 赝 PDF 的匹配系数再次显示出 $z\to 0$ 不是光滑的定域极限。可以定义一个 $\overline{\rm MS}$ 以外的方案以保证该极限光滑，重现与``定域算符对应于守恒流''这一事实相一致的 $z\to 0$ 重正化。这样的一个方案是把所有 $\overline{\rm MS}$ 重正化常数简单地乘以 $C_0(\mu^2 z^2)$，这将给出另一种方案下重正化的\spcorr，以及式~(\ref{eq:ps-c})中相应的、具有光滑 $z\to 0$ 极限的不同匹配系数。这等价于一开始就按照文献~\cite{Radyushkin:2017cyf,Radyushkin:2017lvu}所倡导的方式研究式~(\ref{eq:Qtratio})的比值。下面我们将给出该方案选择的显式结果。这个修改后的方案不应与 $\overline{\rm MS}$ 方案的严格定义相混淆。


由式~(\ref{eq:fac-tq})，与式~(\ref{eq:fac-tq-orig0})定义的准 PDF 相应的匹配系数关系为
\beq \label{eq:quasiC}
C^{(1)}(\xi,\mu/(|y|P^z))= \tilde q'^{(1)}(\xi,\mu/|y|P^z,\epsilon_{\mbox{\tiny IR}})-q^{(1)}(\xi,\epsilon_{\mbox{\tiny IR}})\,.
\eeq
于是利用式~(\ref{eq:qPDFft})，联系准 PDF 与 PDF 的匹配系数是
\begin{align} \label{eq:quasi-c}
C\left(\xi, {\mu\over |y| P^z}\right)
 = &\, \delta\left(1-\xi\right)+{\alpha_sC_F\over 2\pi}\left\{
	\begin{array}{ll}
	\displaystyle \left({1+\xi^2\over 1-\xi}\ln {\xi\over \xi-1} + 1 + {3\over 2\xi}\right)^{[1,\infty]}_{+(1)}- \left({3\over 2\xi}\right)^{[1,\infty]}_{+(\infty)}
    &\, \xi>1
    \\[10pt]
	\displaystyle \left({1+\xi^2\over 1-\xi}\left[
     - \ln{\mu^2\over y^2P_z^2} + \ln\big(4\xi(1-\xi)\big)\right] - {\xi(1+\xi)\over 1-\xi}\right)^{[0,1]}_{+(1)}
    &\, 0<\xi<1
    \\[10pt]
	\displaystyle  \left(-{1+\xi^2\over 1-\xi}\ln {-\xi\over 1-\xi} - 1 + {3\over 2(1-\xi)}\right)^{[-\infty,0]}_{+(1)} - \left({3\over 2(1-\xi)}\right)^{[-\infty,0]}_{+(-\infty)} \quad
    &\, \xi<0
	\end{array}\right.\\[5pt]
	& + {\alpha_sC_F\over 2\pi}\left[\delta(1-\xi) \left( {3\over2}\ln{\mu^2\over 4y^2P_z^2} + {5\over2}\right) + {3\over2}\gamma_E\left({1\over(\xi-1)^2}\delta^+({1\over \xi-1}) + {1\over(1-\xi)^2}\delta^+({1\over 1-\xi})\right)\right]\,.\nn
\end{align}
同样，如预期那样，该结果与红外调节无关。这里 plus 函数项 $\big[ g_1(\xi) \big]_{+(1)}^{[1,\infty]}$ 与 $\big[ g_2(\xi) \big]_{+(1)}^{[-\infty,0]}$ 的被积函数在 $\xi\to \pm \infty$ 时收敛，行为如 $g_i(\xi)\sim 1/\xi^2$。注意，如果我们改用式~(\ref{eq:qPDFren})计算的重正化 $\overline{\rm MS}$ 准 PDF，就会得到在 $\xi=\pm\infty$ 处带不同 $\delta$ 函数的另一个匹配系数 $C$。然而，对任何可积 PDF，这些 $\delta$ 函数对式~(\ref{eq:fac-tq})中的卷积积分没有贡献。例如，为完成与 $1/\xi^2 \delta^+(1/\xi)$ 的卷积，可以利用 $\delta^+\big({1\over \xi}\big)  = \lim_{\beta\to0^+}  \delta\big({1\over \xi} - \beta\big)$，代入因子化公式即给出
\begin{align}
&\lim_{\beta\to0^+}\int {dy\over |y|} {y^2\over x^2}
\, \delta\Big({y\over x} -\beta\Big) f_{u-d}(y)
=\lim_{\beta\to0^+}\beta f_{u-d}(\beta x)\,.
\end{align}
对于 $\infty$ 处的 plus 函数，利用式~(\ref{eq:plusinfinity})与(\ref{eq:plus1})我们有
\begin{align}
\int_{-1}^{+1}\! {dy\over |y|} \left[{1\over (x/y)}\right]_{+(\infty)}^{[1,\infty]} f_{u-d}(y)
&= \lim_{\beta\to 0^+} \int_{-1}^{+1}\! {dy\over |y|}
\bigg[  {\theta(x/y-\beta) \over x/y}
+ {y^2\over x^2}\delta\Bigl({y\over x} -\beta\Bigr)\ln\beta  \bigg] f_{u-d}(y)
\\
&= \int_{-1}^{+1}\! {dy\over x} {y\over |y|} f_{u-d}(y)
+ \lim_{\beta\to 0^+}\beta f_{u-d}(\beta x)\ln\beta
\,. \nn
\end{align}
在最后一行我们丢掉了 $\theta(x/y-\beta)$，因为在小 $y$ 处我们的 PDF 表现为 $f_{u-d}(y)\sim y^{-1+a}$（$0<a<1$）。这也意味着
\begin{align}
\lim_{\beta\to0}\beta f_{u-d}(\beta x) \propto\: & x^{-1+a}\lim_{\beta\to0} \beta^a  =0\,, \\
\lim_{\beta\to0}\beta f_{u-d}(\beta x)\ln\beta \propto\: & x^{-1+a}\lim_{\beta\to0} \beta^a \ln\beta =0\,,
\end{align}
也就是说，匹配系数 $C$ 中在 $\xi=\pm\infty$ 处取值的分布贡献为零。

因此，由式~(\ref{eq:qPDFren})与式~(\ref{eq:qPDFft})的准 PDF 算得的匹配系数实际上相同，把它们代入因子化公式时可以简单地去掉 $\xi=\pm\infty$ 处的所有 $\delta$ 函数：
\begin{align} \label{eq:quasi-matching}
C^{\overline{\rm MS}}\left(\xi, {\mu\over |y| P^z}\right)
= &\, \delta\left(1-\xi\right)+{\alpha_sC_F\over 2\pi}\left\{
\begin{array}{ll}
\displaystyle \left({1+\xi^2\over 1-\xi}\ln {\xi\over \xi-1} + 1 + {3\over 2\xi}\right)^{[1,\infty]}_{+(1)}- {3\over 2\xi}
&\, \xi>1
\\[10pt]
\displaystyle \left({1+\xi^2\over 1-\xi}\left[
- \ln{\mu^2\over y^2P_z^2} + \ln\big(4\xi(1-\xi)\big)\right] - {\xi(1+\xi)\over 1-\xi}\right)^{[0,1]}_{+(1)}
&\, 0<\xi<1
\\[10pt]
\displaystyle  \left(-{1+\xi^2\over 1-\xi}\ln {-\xi\over 1-\xi} - 1 + {3\over 2(1-\xi)}\right)^{[-\infty,0]}_{+(1)} - {3\over 2(1-\xi)}\quad
&\, \xi<0
\end{array}\right.\\[5pt]
& + {\alpha_sC_F\over 2\pi}\delta(1-\xi) \left( {3\over2}\ln{\mu^2\over 4y^2P_z^2} + {5\over2}\right)\,.\nn
\end{align}
把式~(\ref{eq:quasi-matching})用于因子化公式，对任何满足 $\lim_{y\to 0} f(y,\mu) \sim y^{-1+a}$（$a>0$）的 PDF 都有效。
我们还计算了 $\Gamma=\gamma^z$ 情形的匹配系数，它由 $C_{\gamma^z}(\xi, \mu/(|y| P^z)) = C(\xi, \mu/(|y| P^z)) + \Delta C_{\gamma^z}(\xi, \mu/(|y| P^z))$ 给出，其中
\begin{align}
\Delta C_{\gamma^z}(\xi, \mu/(|y| P^z))  = {\alpha_s C_F\over 2\pi} 2(1-\xi)
  \,\theta(\xi)\theta(1-\xi) \,.
\end{align}

注意，我们在 $\overline{\rm MS}$ 下得到的夸克匹配系数不同于文献~\cite{Wang:2017qyg}的结果——后者是一个纯 plus 函数，但其卷积不收敛，就像带横向动量截断的准 PDF 情形一样，参见文献~\cite{Stewart:2017tvs}。

由于重正化赝 PDF 与准 PDF 按定义满足式~(\ref{eq:equiv})，式~(\ref{eq:ps-c}, \ref{eq:quasi-c})给出的 $C\left(\xi, \mu/(|y| P^z)\right)$ 与 ${\cal C}(\alpha,z^2\mu^2)$ 自动满足式~(\ref{eq:ftc-C})中的关系。

此外，如果准 PDF 使用 $\overline{\rm MS}$ 以外的方案——例如通过把 $C_0$ 吸收进 $\overline{\rm MS}$ 重正化常数得到的方案——那么所得匹配系数将是纯 plus 函数从而满足流守恒。从式~(\ref{eq:quasi-c})出发，并用式~(\ref{eq:C0})与式~(\ref{eq:subtlety})，我们得到
\begin{align} \label{eq:quasi-matching-alt}
C^{\rm ratio\!}\left(\xi, {\mu\over |y| P^z}\right)
= &\, \delta\left(1-\xi\right)+{\alpha_sC_F\over 2\pi}\left\{
\begin{array}{ll}
\displaystyle \left({1+\xi^2\over 1-\xi}\ln {\xi\over \xi-1} + 1 - {3\over 2(1-\xi)}\right)^{[1,\infty]}_{+(1)}
&\, \xi>1
\\[10pt]
\displaystyle \left({1+\xi^2\over 1-\xi}\left[
- \ln{\mu^2\over y^2P_z^2} + \ln\big(4\xi(1-\xi)\big)-1\right] +1+ {3\over 2(1-\xi)}\right)^{[0,1]}_{+(1)}
&\, 0<\xi<1
\\[10pt]
\displaystyle  \left(-{1+\xi^2\over 1-\xi}\ln {-\xi\over 1-\xi} - 1 + {3\over 2(1-\xi)}\right)^{[-\infty,0]}_{+(1)}\quad
&\, \xi<0
\end{array}\right.\,,
\end{align}
而 $\Gamma=\gamma^z$ 情形为
\beq
\Delta C^{\rm ratio}_{\gamma^z}(\xi, \mu/(|y| P^z))  = {\alpha_s C_F\over 2\pi} \big[2(1-\xi)\big]^{[0,1]}_{+(1)} \,.
\eeq
在重正化准 PDF 保持流守恒的同时，式~(\ref{eq:quasi-matching-alt})可以用作例如文献~\cite{Alexandrou:2017huk}中 PDF 格点计算的两步匹配程序的输入。
就匹配步骤而言，一个等价的做法是一开始就在 $\overline{\rm MS}$ 方案中研究式~(\ref{eq:Qtratio})给出的比值（如文献~\cite{Radyushkin:2017cyf,Radyushkin:2017lvu}所倡导），并将其匹配到 PDF，这将给出式~(\ref{eq:quasi-matching-alt})。至此我们结束了关于匹配结果以及准 PDF 与赝 PDF 在单圈阶等价性的讨论。

\end{widetext}
